\section{Zusammenfassung und Ausblick}
%==============================================================================

\subsection{Kernbeiträge des META-Moduls}

Das \META-Modul leistet folgende Beiträge zum \BCM:

\begin{enumerate}
    \item \textbf{Axiomatische Fundierung}: 115 explizite Annahmen machen das Modell transparent und kritisierbar
    \item \textbf{Interdisziplinäre Verankerung}: Die vier Kategorien (ONT, EPI, GOV, DYN) verbinden Verhaltensökonomie mit klassischen philosophischen Disziplinen
    \item \textbf{Bias-Katalog}: 37 formalisierte kognitive Biases ermöglichen systematische Berücksichtigung menschlicher Irrationalität
    \item \textbf{Ethische Leitplanken}: Die GOV-Kategorie definiert klare Prinzipien für Interventionsdesign
    \item \textbf{LLM-Guidance}: Die EPI-Axiome informieren die BEATRIX-Parameterschätzung
\end{enumerate}

\subsection{Versionierung}

\begin{table}[h]
\centering
\caption{META-Modul Versionshistorie}
\begin{tabular}{@{}llll@{}}
\toprule
\textbf{Version} & \textbf{Datum} & \textbf{Axiome} & \textbf{Änderungen} \\
\midrule
1.0 & 2023 & 100 & Initiale Version \\
2.0 & 2024 & 100 & BCM 2.2 Integration \\
2.1 & Dez 2024 & 101 & MA-ONT-26: 5-Ebenen-Hierarchie \\
2.2 & Dez 2024 & 114 & +12 Biases (EPI-26 bis EPI-37), +Evidenzhierarchie (EPI-38) \\
2.3 & Dez 2024 & 115 & +End-to-End Integration (INF-08) \\
\bottomrule
\end{tabular}
\end{table}

\subsection{Zukünftige Erweiterungen}

Geplante Erweiterungen für \META{} 3.0:
\begin{itemize}
    \item Integration von Affekt-Dynamik-Axiomen
    \item Formalisierung von Sozialkapital-Konzepten
    \item Erweiterung der Governance-Kategorie um KI-spezifische Ethik
    \item Cross-kulturelle Validierung der Bias-Universalität
\end{itemize}

%==============================================================================
% References
%==============================================================================

\section*{Literaturverzeichnis}

\subsection*{Grundlagenwerke Verhaltensökonomie}

\begin{itemize}[leftmargin=*]
    \item Kahneman, D. (2011). \textit{Thinking, Fast and Slow}. Farrar, Straus and Giroux.
    \item Kahneman, D., \& Tversky, A. (1979). Prospect theory: An analysis of decision under risk. \textit{Econometrica}, 47(2), 263-291.
    \item Thaler, R. H. (2015). \textit{Misbehaving: The Making of Behavioral Economics}. W. W. Norton.
    \item Thaler, R. H., \& Sunstein, C. R. (2008). \textit{Nudge: Improving Decisions About Health, Wealth, and Happiness}. Yale University Press.
    \item Ariely, D. (2008). \textit{Predictably Irrational: The Hidden Forces That Shape Our Decisions}. HarperCollins.
    \item Camerer, C. F., Loewenstein, G., \& Rabin, M. (Eds.). (2004). \textit{Advances in Behavioral Economics}. Princeton University Press.
    \item Gilovich, T., Griffin, D., \& Kahneman, D. (Eds.). (2002). \textit{Heuristics and Biases: The Psychology of Intuitive Judgment}. Cambridge University Press.
\end{itemize}

\subsection*{Bounded Rationality \& Heuristiken (MA-EPI-01 bis MA-EPI-06)}

\begin{itemize}[leftmargin=*]
    \item Simon, H. A. (1955). A behavioral model of rational choice. \textit{Quarterly Journal of Economics}, 69(1), 99-118.
    \item Simon, H. A. (1957). \textit{Models of Man: Social and Rational}. Wiley.
    \item Gigerenzer, G. (2007). \textit{Gut Feelings: The Intelligence of the Unconscious}. Viking.
    \item Gigerenzer, G., \& Selten, R. (Eds.). (2001). \textit{Bounded Rationality: The Adaptive Toolbox}. MIT Press.
    \item Gigerenzer, G., Todd, P. M., \& ABC Research Group. (1999). \textit{Simple Heuristics That Make Us Smart}. Oxford University Press.
    \item Tversky, A., \& Kahneman, D. (1974). Judgment under uncertainty: Heuristics and biases. \textit{Science}, 185(4157), 1124-1131.
    \item Stanovich, K. E., \& West, R. F. (2000). Individual differences in reasoning: Implications for the rationality debate. \textit{Behavioral and Brain Sciences}, 23(5), 645-665.
    \item Evans, J. S. B. T. (2008). Dual-processing accounts of reasoning, judgment, and social cognition. \textit{Annual Review of Psychology}, 59, 255-278.
\end{itemize}

\subsection*{Kognitive Biases -- Klassische Heuristiken (MA-EPI-12 bis MA-EPI-18)}

\begin{itemize}[leftmargin=*]
    \item Tversky, A., \& Kahneman, D. (1973). Availability: A heuristic for judging frequency and probability. \textit{Cognitive Psychology}, 5(2), 207-232.
    \item Tversky, A., \& Kahneman, D. (1981). The framing of decisions and the psychology of choice. \textit{Science}, 211(4481), 453-458.
    \item Tversky, A., \& Kahneman, D. (1992). Advances in prospect theory: Cumulative representation of uncertainty. \textit{Journal of Risk and Uncertainty}, 5(4), 297-323.
    \item Kahneman, D., Knetsch, J. L., \& Thaler, R. H. (1991). Anomalies: The endowment effect, loss aversion, and status quo bias. \textit{Journal of Economic Perspectives}, 5(1), 193-206.
    \item Samuelson, W., \& Zeckhauser, R. (1988). Status quo bias in decision making. \textit{Journal of Risk and Uncertainty}, 1(1), 7-59.
    \item Strack, F., \& Mussweiler, T. (1997). Explaining the enigmatic anchoring effect: Mechanisms of selective accessibility. \textit{Journal of Personality and Social Psychology}, 73(3), 437-446.
    \item Epley, N., \& Gilovich, T. (2006). The anchoring-and-adjustment heuristic: Why the adjustments are insufficient. \textit{Psychological Science}, 17(4), 311-318.
    \item Moore, D. A., \& Healy, P. J. (2008). The trouble with overconfidence. \textit{Psychological Review}, 115(2), 502-517.
    \item Fischhoff, B., Slovic, P., \& Lichtenstein, S. (1977). Knowing with certainty: The appropriateness of extreme confidence. \textit{Journal of Experimental Psychology: Human Perception and Performance}, 3(4), 552-564.
\end{itemize}

\subsection*{Kognitive Biases -- Erweiterte Liste (MA-EPI-26 bis MA-EPI-37)}

\begin{itemize}[leftmargin=*]
    \item Kahneman, D., \& Tversky, A. (1973). On the psychology of prediction. \textit{Psychological Review}, 80(4), 237-251. [Base Rate Neglect]
    \item Tversky, A., \& Kahneman, D. (1972). Subjective probability: A judgment of representativeness. \textit{Cognitive Psychology}, 3(3), 430-454. [Representativeness]
    \item Arkes, H. R., \& Blumer, C. (1985). The psychology of sunk cost. \textit{Organizational Behavior and Human Decision Processes}, 35(1), 124-140. [Sunk Cost]
    \item Thaler, R. (1980). Toward a positive theory of consumer choice. \textit{Journal of Economic Behavior \& Organization}, 1(1), 39-60. [Endowment Effect]
    \item Kahneman, D., Knetsch, J. L., \& Thaler, R. H. (1990). Experimental tests of the endowment effect and the Coase theorem. \textit{Journal of Political Economy}, 98(6), 1325-1348.
    \item Weinstein, N. D. (1980). Unrealistic optimism about future life events. \textit{Journal of Personality and Social Psychology}, 39(5), 806-820. [Optimism Bias]
    \item Sharot, T. (2011). The optimism bias. \textit{Current Biology}, 21(23), R941-R945.
    \item Miller, D. T., \& Ross, M. (1975). Self-serving biases in the attribution of causality. \textit{Psychological Bulletin}, 82(2), 213-225. [Self-Serving Bias]
    \item Crowne, D. P., \& Marlowe, D. (1960). A new scale of social desirability. \textit{Journal of Consulting Psychology}, 24(4), 349-354. [Social Desirability]
    \item Paulhus, D. L. (1984). Two-component models of socially desirable responding. \textit{Journal of Personality and Social Psychology}, 46(3), 598-609.
    \item Spranca, M., Minsk, E., \& Baron, J. (1991). Omission and commission in judgment and choice. \textit{Journal of Experimental Social Psychology}, 27(1), 76-105. [Omission Bias]
    \item Ritov, I., \& Baron, J. (1990). Reluctance to vaccinate: Omission bias and ambiguity. \textit{Journal of Behavioral Decision Making}, 3(4), 263-277.
    \item Merritt, A. C., Effron, D. A., \& Monin, B. (2010). Moral self-licensing: When being good frees us to be bad. \textit{Social and Personality Psychology Compass}, 4(5), 344-357. [Moral Licensing]
    \item Sachdeva, S., Iliev, R., \& Medin, D. L. (2009). Sinning saints and saintly sinners: The paradox of moral self-regulation. \textit{Psychological Science}, 20(4), 523-528.
    \item Small, D. A., \& Loewenstein, G. (2003). Helping a victim or helping the victim: Altruism and identifiability. \textit{Journal of Risk and Uncertainty}, 26(1), 5-16. [Identifiable Victim]
    \item Jenni, K., \& Loewenstein, G. (1997). Explaining the identifiable victim effect. \textit{Journal of Risk and Uncertainty}, 14(3), 235-257.
    \item Kahneman, D. (1986). Comments on the contingent valuation method. In R. G. Cummings et al. (Eds.), \textit{Valuing Environmental Goods}. Rowman \& Allanheld. [Scope Insensitivity]
    \item Desvousges, W. H., et al. (1993). Measuring nonuse damages using contingent valuation. \textit{Journal of Environmental Economics and Management}, 23(3), 177-194.
    \item Slovic, P., Finucane, M., Peters, E., \& MacGregor, D. G. (2002). The affect heuristic. In T. Gilovich et al. (Eds.), \textit{Heuristics and Biases} (pp. 397-420). Cambridge University Press. [Affect Heuristic]
    \item Slovic, P., Finucane, M. L., Peters, E., \& MacGregor, D. G. (2007). The affect heuristic. \textit{European Journal of Operational Research}, 177(3), 1333-1352.
    \item Nickerson, R. S. (1998). Confirmation bias: A ubiquitous phenomenon in many guises. \textit{Review of General Psychology}, 2(2), 175-220. [Confirmation Bias]
\end{itemize}

\subsection*{Identitätsökonomie \& Soziale Präferenzen (MA-ONT-09 bis MA-ONT-11)}

\begin{itemize}[leftmargin=*]
    \item Akerlof, G. A., \& Kranton, R. E. (2000). Economics and identity. \textit{Quarterly Journal of Economics}, 115(3), 715-753.
    \item Akerlof, G. A., \& Kranton, R. E. (2010). \textit{Identity Economics: How Our Identities Shape Our Work, Wages, and Well-Being}. Princeton University Press.
    \item Bénabou, R., \& Tirole, J. (2011). Identity, morals, and taboos: Beliefs as assets. \textit{Quarterly Journal of Economics}, 126(2), 805-855.
    \item Fehr, E., \& Schmidt, K. M. (1999). A theory of fairness, competition, and cooperation. \textit{Quarterly Journal of Economics}, 114(3), 817-868.
    \item Fehr, E., \& Fischbacher, U. (2002). Why social preferences matter -- the impact of non-selfish motives on competition, cooperation and incentives. \textit{Economic Journal}, 112(478), C1-C33.
    \item Fischbacher, U., Gächter, S., \& Fehr, E. (2001). Are people conditionally cooperative? Evidence from a public goods experiment. \textit{Economics Letters}, 71(3), 397-404.
    \item Bolton, G. E., \& Ockenfels, A. (2000). ERC: A theory of equity, reciprocity, and competition. \textit{American Economic Review}, 90(1), 166-193.
    \item Charness, G., \& Rabin, M. (2002). Understanding social preferences with simple tests. \textit{Quarterly Journal of Economics}, 117(3), 817-869.
    \item Tajfel, H. (1978). \textit{Differentiation Between Social Groups}. Academic Press.
    \item Tajfel, H., \& Turner, J. C. (1979). An integrative theory of intergroup conflict. In W. G. Austin \& S. Worchel (Eds.), \textit{The Social Psychology of Intergroup Relations} (pp. 33-47). Brooks/Cole.
\end{itemize}

\subsection*{Soziale Normen \& Kooperation (MA-ONT-11, MA-DYN-10 bis MA-DYN-12)}

\begin{itemize}[leftmargin=*]
    \item Bicchieri, C. (2006). \textit{The Grammar of Society: The Nature and Dynamics of Social Norms}. Cambridge University Press.
    \item Bicchieri, C. (2017). \textit{Norms in the Wild: How to Diagnose, Measure, and Change Social Norms}. Oxford University Press.
    \item Cialdini, R. B., Reno, R. R., \& Kallgren, C. A. (1990). A focus theory of normative conduct: Recycling the concept of norms. \textit{Journal of Personality and Social Psychology}, 58(6), 1015-1026.
    \item Cialdini, R. B., \& Goldstein, N. J. (2004). Social influence: Compliance and conformity. \textit{Annual Review of Psychology}, 55, 591-621.
    \item Centola, D. (2018). \textit{How Behavior Spreads: The Science of Complex Contagions}. Princeton University Press.
    \item Centola, D., \& Macy, M. (2007). Complex contagions and the weakness of long ties. \textit{American Journal of Sociology}, 113(3), 702-734.
    \item Granovetter, M. (1978). Threshold models of collective behavior. \textit{American Journal of Sociology}, 83(6), 1420-1443.
    \item Young, H. P. (2015). The evolution of social norms. \textit{Annual Review of Economics}, 7, 359-387.
    \item Ostrom, E. (1990). \textit{Governing the Commons: The Evolution of Institutions for Collective Action}. Cambridge University Press.
\end{itemize}

\subsection*{Temporale Entscheidungen \& Discounting (MA-DYN-13 bis MA-DYN-16)}

\begin{itemize}[leftmargin=*]
    \item Frederick, S., Loewenstein, G., \& O'Donoghue, T. (2002). Time discounting and time preference: A critical review. \textit{Journal of Economic Literature}, 40(2), 351-401.
    \item Laibson, D. (1997). Golden eggs and hyperbolic discounting. \textit{Quarterly Journal of Economics}, 112(2), 443-478.
    \item O'Donoghue, T., \& Rabin, M. (1999). Doing it now or later. \textit{American Economic Review}, 89(1), 103-124.
    \item Strotz, R. H. (1955). Myopia and inconsistency in dynamic utility maximization. \textit{Review of Economic Studies}, 23(3), 165-180.
    \item Loewenstein, G. (1996). Out of control: Visceral influences on behavior. \textit{Organizational Behavior and Human Decision Processes}, 65(3), 272-292.
    \item Loewenstein, G., \& Prelec, D. (1992). Anomalies in intertemporal choice: Evidence and an interpretation. \textit{Quarterly Journal of Economics}, 107(2), 573-597.
    \item Loewenstein, G., O'Donoghue, T., \& Rabin, M. (2003). Projection bias in predicting future utility. \textit{Quarterly Journal of Economics}, 118(4), 1209-1248.
    \item Van Boven, L., \& Loewenstein, G. (2003). Social projection of transient drive states. \textit{Personality and Social Psychology Bulletin}, 29(9), 1159-1168.
\end{itemize}

\subsection*{Gewohnheiten \& Selbstkontrolle (MA-DYN-07 bis MA-DYN-09, MA-DYN-17 bis MA-DYN-18)}

\begin{itemize}[leftmargin=*]
    \item Duhigg, C. (2012). \textit{The Power of Habit: Why We Do What We Do in Life and Business}. Random House.
    \item Wood, W., \& Rünger, D. (2016). Psychology of habit. \textit{Annual Review of Psychology}, 67, 289-314.
    \item Neal, D. T., Wood, W., \& Quinn, J. M. (2006). Habits -- A repeat performance. \textit{Current Directions in Psychological Science}, 15(4), 198-202.
    \item Baumeister, R. F., Bratslavsky, E., Muraven, M., \& Tice, D. M. (1998). Ego depletion: Is the active self a limited resource? \textit{Journal of Personality and Social Psychology}, 74(5), 1252-1265.
    \item Baumeister, R. F., Vohs, K. D., \& Tice, D. M. (2007). The strength model of self-control. \textit{Current Directions in Psychological Science}, 16(6), 351-355.
    \item Gollwitzer, P. M. (1999). Implementation intentions: Strong effects of simple plans. \textit{American Psychologist}, 54(7), 493-503.
    \item Gollwitzer, P. M., \& Sheeran, P. (2006). Implementation intentions and goal achievement: A meta-analysis. \textit{Advances in Experimental Social Psychology}, 38, 69-119.
    \item Bryan, G., Karlan, D., \& Nelson, S. (2010). Commitment devices. \textit{Annual Review of Economics}, 2(1), 671-698.
\end{itemize}

\subsection*{Hedonische Adaptation \& Wohlbefinden (MA-DYN-19 bis MA-DYN-21)}

\begin{itemize}[leftmargin=*]
    \item Brickman, P., \& Campbell, D. T. (1971). Hedonic relativism and planning the good society. In M. H. Appley (Ed.), \textit{Adaptation-Level Theory} (pp. 287-305). Academic Press.
    \item Brickman, P., Coates, D., \& Janoff-Bulman, R. (1978). Lottery winners and accident victims: Is happiness relative? \textit{Journal of Personality and Social Psychology}, 36(8), 917-927.
    \item Diener, E., Lucas, R. E., \& Scollon, C. N. (2006). Beyond the hedonic treadmill: Revising the adaptation theory of well-being. \textit{American Psychologist}, 61(4), 305-314.
    \item Kahneman, D., Wakker, P. P., \& Sarin, R. (1997). Back to Bentham? Explorations of experienced utility. \textit{Quarterly Journal of Economics}, 112(2), 375-406.
    \item Kahneman, D., Fredrickson, B. L., Schreiber, C. A., \& Redelmeier, D. A. (1993). When more pain is preferred to less: Adding a better end. \textit{Psychological Science}, 4(6), 401-405. [Peak-End Rule]
    \item Redelmeier, D. A., \& Kahneman, D. (1996). Patients' memories of painful medical treatments: Real-time and retrospective evaluations. \textit{Pain}, 66(1), 3-8.
\end{itemize}

\subsection*{Nudging \& Interventionsdesign (MA-GOV-07 bis MA-GOV-10)}

\begin{itemize}[leftmargin=*]
    \item Thaler, R. H., \& Sunstein, C. R. (2008). \textit{Nudge: Improving Decisions About Health, Wealth, and Happiness}. Yale University Press.
    \item Sunstein, C. R. (2014). \textit{Why Nudge? The Politics of Libertarian Paternalism}. Yale University Press.
    \item Johnson, E. J., \& Goldstein, D. (2003). Do defaults save lives? \textit{Science}, 302(5649), 1338-1339.
    \item Madrian, B. C., \& Shea, D. F. (2001). The power of suggestion: Inertia in 401(k) participation and savings behavior. \textit{Quarterly Journal of Economics}, 116(4), 1149-1187.
    \item Benartzi, S., \& Thaler, R. H. (2013). Behavioral economics and the retirement savings crisis. \textit{Science}, 339(6124), 1152-1153.
    \item Hertwig, R., \& Grüne-Yanoff, T. (2017). Nudging and boosting: Steering or empowering good decisions. \textit{Perspectives on Psychological Science}, 12(6), 973-986.
    \item Grüne-Yanoff, T., \& Hertwig, R. (2016). Nudge versus boost: How coherent are policy and theory? \textit{Minds and Machines}, 26(1-2), 149-183.
    \item Hallsworth, M., \& Kirkman, E. (2020). \textit{Behavioral Insights}. MIT Press.
\end{itemize}

\subsection*{Ethik \& Governance von Interventionen (MA-GOV-03 bis MA-GOV-06, MA-GOV-11 bis MA-GOV-18)}

\begin{itemize}[leftmargin=*]
    \item Sunstein, C. R. (2016). \textit{The Ethics of Influence: Government in the Age of Behavioral Science}. Cambridge University Press.
    \item Hausman, D. M., \& Welch, B. (2010). Debate: To nudge or not to nudge. \textit{Journal of Political Philosophy}, 18(1), 123-136.
    \item Hansen, P. G., \& Jespersen, A. M. (2013). Nudge and the manipulation of choice. \textit{European Journal of Risk Regulation}, 4(1), 3-28.
    \item Bovens, L. (2009). The ethics of nudge. In T. Grüne-Yanoff \& S. O. Hansson (Eds.), \textit{Preference Change} (pp. 207-219). Springer.
    \item Frey, B. S., \& Jegen, R. (2001). Motivation crowding theory. \textit{Journal of Economic Surveys}, 15(5), 589-611. [Crowding Out]
    \item Deci, E. L., Koestner, R., \& Ryan, R. M. (1999). A meta-analytic review of experiments examining the effects of extrinsic rewards on intrinsic motivation. \textit{Psychological Bulletin}, 125(6), 627-668.
    \item Gneezy, U., Meier, S., \& Rey-Biel, P. (2011). When and why incentives (don't) work to modify behavior. \textit{Journal of Economic Perspectives}, 25(4), 191-210.
    \item Bowles, S., \& Polania-Reyes, S. (2012). Economic incentives and social preferences: Substitutes or complements? \textit{Journal of Economic Literature}, 50(2), 368-425.
\end{itemize}

\subsection*{Dynamische Systeme \& Komplexität (MA-DYN-04 bis MA-DYN-06)}

\begin{itemize}[leftmargin=*]
    \item Arthur, W. B. (1994). \textit{Increasing Returns and Path Dependence in the Economy}. University of Michigan Press.
    \item David, P. A. (1985). Clio and the economics of QWERTY. \textit{American Economic Review}, 75(2), 332-337.
    \item Gladwell, M. (2000). \textit{The Tipping Point: How Little Things Can Make a Big Difference}. Little, Brown.
    \item Scheffer, M. (2009). \textit{Critical Transitions in Nature and Society}. Princeton University Press.
    \item Page, S. E. (2006). Path dependence. \textit{Quarterly Journal of Political Science}, 1(1), 87-115.
\end{itemize}

\subsection*{Methodologie \& Evaluation (MA-GOV-20 bis MA-GOV-24)}

\begin{itemize}[leftmargin=*]
    \item Angrist, J. D., \& Pischke, J. S. (2009). \textit{Mostly Harmless Econometrics: An Empiricist's Companion}. Princeton University Press.
    \item Duflo, E., Glennerster, R., \& Kremer, M. (2007). Using randomization in development economics research: A toolkit. \textit{Handbook of Development Economics}, 4, 3895-3962.
    \item List, J. A. (2011). Why economists should conduct field experiments and 14 tips for pulling one off. \textit{Journal of Economic Perspectives}, 25(3), 3-16.
    \item DellaVigna, S., \& Linos, E. (2022). RCTs to scale: Comprehensive evidence from two nudge units. \textit{Econometrica}, 90(1), 81-116.
    \item Banerjee, A. V., \& Duflo, E. (2009). The experimental approach to development economics. \textit{Annual Review of Economics}, 1, 151-178.
\end{itemize}

\subsection*{Kulturelle \& Institutionelle Ebenen (MA-ONT-26)}

\begin{itemize}[leftmargin=*]
    \item Hofstede, G. (2001). \textit{Culture's Consequences: Comparing Values, Behaviors, Institutions and Organizations Across Nations} (2nd ed.). Sage.
    \item Henrich, J., Heine, S. J., \& Norenzayan, A. (2010). The weirdest people in the world? \textit{Behavioral and Brain Sciences}, 33(2-3), 61-83.
    \item North, D. C. (1990). \textit{Institutions, Institutional Change and Economic Performance}. Cambridge University Press.
    \item Acemoglu, D., \& Robinson, J. A. (2012). \textit{Why Nations Fail: The Origins of Power, Prosperity, and Poverty}. Crown Business.
    \item Alesina, A., \& Giuliano, P. (2015). Culture and institutions. \textit{Journal of Economic Literature}, 53(4), 898-944.
    \item Gelfand, M. J. (2018). \textit{Rule Makers, Rule Breakers: How Tight and Loose Cultures Wire Our World}. Scribner.
\end{itemize}

%==============================================================================
% Appendix
%==============================================================================

\appendix

